\section{Clustering}
\begin{frame}{\hspace{3cm}}
\centering

\vspace{1.5cm}

\begin{beamercolorbox}[sep=20pt,center,rounded=true,shadow=true]{title}
    {\fontsize{20}{24}\selectfont \textbf{Clustering}}
\end{beamercolorbox}

\vspace{0.8cm}


\vfill

\end{frame}
%------------------------------------------------
\begin{frame}{What does clustering stand for?}

\begin{block}{Clustering}
The clustering process aims at grouping together points having high similarity, 
placing those with large dissimilarity into different groups.
\end{block}

\vspace{0.3cm}

\begin{columns}

\column{0.45\textwidth}
\textbf{Various Types of Clustering}
\begin{itemize}
    \item Traditional Approaches
    \begin{itemize}
    \item Partition-based
    \item Distance-based
    \item Grid-based
     \end{itemize}
    \item Model-based
    \item Density-based
\end{itemize}

\end{columns}

\end{frame}
\begin{frame}{Model-Based Clustering}

In model-based clustering, a finite mixture of parametric distributions is introduced. More precisely, the random vector $X$ is considered to be a mixture of $J$ sub-populations.

\vspace{0.3cm}

\begin{block}{Density Model}
\[
f(x)= \sum_{j=1}^{J} \pi_j f_j(x)
\]
\end{block}

where $f_{j}$ denotes the density of the $j$-th sub-population and $\pi_{j}$ represents its relative weight.

\vspace{0.3cm}

In this framework, a cluster is associated with each component $f_{j}$, and any observation $x_{0}$ is allocated to the cluster with the maximal density $f_{j}(x_{0})$ among the $J$ components.

\end{frame}
\begin{frame}{Model-Based Clustering: Drawbacks}

\begin{alertblock}{Limitations}
\begin{itemize}
    \item To render the estimation problem identifiable, some restrictions must be placed on the components.
    \item The number of components should be determined in advance.
    \item If the cluster shapes do not match those of the components $f_j$, the model-based approach may face difficulties.
\end{itemize}
\end{alertblock}

\end{frame}
\begin{frame}{Density-Based Clustering}

The mechanism behind density-based clustering is to associate groups with \textbf{connected regions of the density}.

\vspace{0.3cm}

\begin{figure}
    \centering
    \includegraphics[width=0.5\linewidth]{WhatsApp Image 2026-04-10 at 15.40.08.jpeg}
\end{figure}

\vspace{0.3cm}


\textbf{Key Idea}
\begin{itemize}
    \item Clusters correspond to \textbf{high-density regions (modes)}
    \item Separation occurs in \textbf{low-density areas}
\end{itemize}

\end{frame}
\begin{frame}{Kernel Density Estimation (KDE)}

Suppose that $X_1, \ldots, X_n$ is a random sample from a density $f$.

\vspace{0.3cm}

\begin{exampleblock}{KDE Formula}
\[
\hat{f}(x) = \frac{1}{n} \sum_{i=1}^{n} \frac{1}{h} \, \kappa\!\left(\frac{x - X_i}{h}\right)
\]
\end{exampleblock}

\vspace{0.3cm}

\textbf{Components}
\begin{itemize}
    \item $\kappa(\cdot)$: kernel function
    \item $h$: bandwidth (controls smoothness)
\end{itemize}

\end{frame}
\begin{frame}{The Non-Parametric pdfCluster}
\begin{itemize}
    \item \textbf{Azzalini and Torelli (2007):} Proposed a \textit{density-based clustering} procedure based on \textit{Kernel Density Estimation (KDE)}. The method exploits \textit{spatial tessellation} to generalize the concept of \textit{adjacency} in multidimensional spaces and to identify \textit{connected level sets}. Clusters are defined as \textit{maximally connected components} of regions where the estimated density exceeds a given threshold.

    \item \textbf{Limitation:} The algorithm is computationally feasible only for \textit{low to moderate dimensional data} (typically up to dimension 6).

    \item \textbf{Menardi and Azzalini (2014):} Extended the methodology by providing an improved approach for detecting \textit{connected high-density regions} in \textit{higher-dimensional spaces}.

    \item \textbf{Implementation:} Both methods are unified into a single algorithm and implemented in the R package \texttt{pdfCluster} (Azzalini and Menardi, 2014).
\end{itemize}
    
\end{frame}

\begin{frame}{Types of Clustering Approach}
\begin{figure}
    \centering
    \includegraphics[width=0.71\linewidth]{Screenshot 2026-04-11 005427.png}
    \caption{Enter Caption}
    \label{fig:placeholder}
\end{figure}
\end{frame}
\begin{frame}{Why Asymmetric Kernel?}
    \begin{itemize}
    \item Classical kernel methods are based on symmetric densities, typically centered around zero, as in the case of the normal kernel.

    \item However, a deficiency arises when such kernels are used to estimate densities whose supports are bounded \textit{[Marron and Ruppert (1994)]}.

    \item On the other hand, asymmetric kernel methods never assign weight outside the support of the density. As a result, they effectively address the boundary bias problem \textit{[Saulo et al., 2013]}.

    \item Density estimators in moderately-sized samples tend to inherit the salient properties of their kernels. Thus, if a density is suspected to be highly skewed, it is more appropriate to use asymmetric kernels rather than symmetric ones \textit{[Abadir and Lawford (2004)]}.
\end{itemize}

\end{frame}
\begin{frame}{Proposed Approach: The Semi-Parametric pdfcluster}
\begin{itemize}
    \item Our main goal in this work is to improve the density estimation component of the pdfCluster algorithm \textit{[Azzalini and Menardi (2014)]} by utilizing skew-symmetric kernels instead of the normal and $t_7$ kernels.

    \item As the terminology suggests, skew-symmetric (or skew-modulated) distributions are capable of generating both symmetric and asymmetric densities; see, for example, \textit{[Salehi and Azzalini (2018)]}.

    \item Among the many distributions in this family, the skew-normal and skew-$t$ distributions—being the most well-known members—are sufficiently flexible to be employed as kernels in kernel density estimation (KDE).
\end{itemize}
    
\end{frame}
%-------------------------
